This section situates our work in the literature on GPU kernel DSLs (\cref{sec:rel:dsls}),
convolution optimization algorithms (\cref{sec:rel:conv}),
performance analysis tools (\cref{sec:rel:profiling}),
and GPU benchmarking frameworks (\cref{sec:rel:bench}).
Our study is the first to show that prevailing DSL and LLM-generated-kernel benchmarks admit
performance-poor kernels, to characterize the resulting hidden gap with a hardware-grounded
root-cause taxonomy, and to distill evaluation heuristics that flag poor kernels absent a
comprehensive benchmark.

\subsection{GPU Kernel DSLs}
\label{sec:rel:dsls}

Halide~\cite{ragan-kelley2013halide} introduced the influential separation of
algorithm specification from scheduling,
enabling systematic search over loop transformations (tiling, vectorization, parallelism)
without modifying the functional description.
This principle informs both TVM~\cite{chen2018tvm} and TileLang~\cite{wang2025tilelang}.

Triton~\cite{tillet2019triton} departed from the explicit-schedule model
in favor of an implicit tile abstraction in which the compiler infers
shared memory layouts, synchronization, and pipelining.
Its integration into PyTorch~\cite{li2023torchcompile} has made it the
most widely deployed GPU DSL.
ThunderKittens~\cite{spector2024thunderkittens} targets the opposite extreme:
a thin C++ header library that exposes warp-level tile operations directly,
achieving state-of-the-art attention throughput on H100 through explicit warp specialization.
Our study is complementary: we characterize performance at the DSL level
rather than at the hand-optimized C++ level.

TVM~\cite{chen2018tvm} and its auto-scheduler Ansor~\cite{zheng2020ansor}
demonstrate that hierarchical program search
substantially outperforms template-guided auto-tuning for irregular operator shapes,
including convolution.
This is consistent with RC2 in our analysis (\cref{sec:analysis}):
the convolution search space is irregularly shaped relative to GEMM,
and static configuration lists are insufficient.
MLIR-based code generation~\cite{bik2022mlir,vasilache2022composable}
offers a compiler-infrastructure approach to the same problem,
generating near-peak-performance GEMM and fused-attention code
through parametric Linalg transformations.

\subsection{Convolution Optimization}
\label{sec:rel:conv}

Zhou~\etal~\cite{zhou2021implicit} provide a thorough characterization of
the implicit GEMM convolution algorithm on GPU Tensor Cores,
identifying the advantages of NHWC layout and on-the-fly tile formation
over explicit im2col.
Their analysis motivates the baseline construction in \cref{sec:meth:baselines}.

Winograd convolution~\cite{lavin2016fast} reduces arithmetic by up to $4\times$
for $3 \times 3$ filters at the cost of additional data transformation passes.
In our measurements (\cref{sec:mitig:conv}), the wall-clock benefit attributable specifically to Winograd selection is much smaller ($\approx$2--3\%); the arithmetic reduction does not translate proportionally because cuDNN's general implicit-GEMM tuning dominates the gap.
cuDNN's selection of Winograd for appropriate shapes is a key component
of its performance lead, as analyzed in \cref{sec:analysis}.
Adding Winograd support as a first-class lowering pass within a DSL like Triton
remains an open engineering problem; our analysis identifies it as future work (\cref{sec:mitig:conv}).

\subsection{Performance Analysis and Profiling}
\label{sec:rel:profiling}

TritonForge~\cite{li2025tritonforge} proposes a profiling-guided LLM loop
for automated Triton kernel optimization,
using Nsight Compute metrics to identify bottlenecks.
Their finding that memory coalescing failures and low occupancy
account for the majority of kernel-level inefficiencies
is consistent with our RC1 and RC3 findings.
The difference is scope:
TritonForge focuses on automation of the fix-generation step,
while our study focuses on systematic characterization across a kernel taxonomy.

Empirical performance studies of GPU libraries have examined cuDNN~\cite{chetlur2014cudnn},
cuBLAS~\cite{cublas2024}, and CUTLASS~\cite{cutlass2023} extensively,
but performance comparisons \emph{between} DSL and library implementations
at the scale and granularity of our study are absent from the literature.

\subsection{Benchmarking GPU Software}
\label{sec:rel:bench}

The benchmarks closest to our setting evaluate DSL and LLM-generated kernels.
TritonBench~\cite{li2025tritonbench} evaluates Triton kernel generation by LLMs,
measuring functional correctness and GPU efficiency as a fraction of hardware peak---
a roofline-anchored metric that predates and motivates the anchor we adopt in \cref{sec:mitig:heuristics}, and which we build on rather than introduce.
KernelBench~\cite{kernelbench} is the de-facto benchmark for LLM kernel generation,
but it gates on correctness alone and reports speedup only as an outcome;
its single reward-hacking guard flags only implausibly fast kernels, a weakness automated generators have exploited.
Both evaluate only a narrow slice of the (DSL, data type, shape, hardware) space.
Our work differs in object and claim: we do not propose another benchmark, but show that the prevailing ones admit performance-poor kernels (\cref{sec:evaluation}), characterize the resulting hidden gap with hardware counters, and distill evaluation heuristics and optimization patterns that guide development absent a comprehensive benchmark.

MLPerf Inference~\cite{reddi2020mlperf} benchmarks end-to-end inference throughput
but treats the computation stack as a black box.
Our study operates at the kernel level
to attribute performance differences to specific compiler behaviors.
